\chapter{Retrieval-Augmented Generation (RAG)}

\section{Theoretical Foundations}
    \subsection{Background}
        \subsubsection{Large Language Model}
         A large language model (LLM) is a type of machine learning model trained on massive datasets using deep neural networks. It is primarily designed for natural language processing (NLP) tasks. The field has experienced a major breakthrough with the introduction of the transformer architecture \cite{transformers}, which relies on the attention mechanism to capture dependencies between words, regardless of their distance in a sequence. Based on this architecture, most existing LLMs fall into one of three categories: encoder-only, decoder-only, or encoder-decoder models.
     \begin{itemize}
        \item \textbf{Encoder-only models} (like BERT) process input in both left-to-right and right-to-left directions simultaneously, enabling them to capture bidirectional context. These models encode the input into a high-dimensional representation and are particularly useful for understanding tasks such as sentiment analysis, text classification, and named entity recognition.
        \item \textbf{Decoder-only models} (like GPT) process input in a left-to-right (autoregressive) manner, predicting the next token based on the sequence of previous tokens. These models excel at generative tasks such as text completion, story generation, and dialogue modeling.
        \item \textbf{Encoder-decoder models} (like T5 or BART) combine both components: the encoder processes the input sequence into a contextual representation, and the decoder generates an output sequence based on this representation. This architecture is well-suited for sequence-to-sequence tasks such as machine translation, summarization, and question answering.
    \end{itemize}
        \subsubsection{Prompt Engineering}

Large language models (LLMs) have demonstrated an impressive ability to perform a wide variety of tasks with minimal or no task-specific training. This versatility gave rise to the field of \textit{prompt engineering}, which involves crafting effective input prompts to steer the model toward producing accurate and relevant outputs for downstream tasks. In the early stages of prompt engineering, users often relied on manually handcrafted templates tailored to each specific task, such as classification, summarization, or question answering. While this method showed promising results, it also presented several limitations: it was highly sensitive to the wording of the prompt, required significant human expertise, and often lacked robustness when applied across different domains or inputs.

To overcome the limitations of handcrafted prompts, a more dynamic technique known as \textit{in-context learning} was introduced. This approach consists in providing the model with a small number of input-output examples related to the task directly within the prompt. Instead of fine-tuning the model's parameters, we guide its behavior by demonstrating what the task looks like. For example, to teach a model how to translate, one might provide a few sentences in English along with their French translations. In-context learning leverages the model's ability to generalize from these demonstrations and generate outputs consistent with the observed examples. This method has been shown to significantly improve performance on many NLP tasks, particularly when fine-tuning is infeasible due to time or resource constraints.

However, in-context learning also has its own limitations. Its effectiveness depends heavily on the quality, relevance, and ordering of the examples presented in the prompt. Furthermore, the model remains confined to the knowledge encoded during its pretraining phase. This means it cannot answer questions about events or facts that occurred after its training cutoff. For instance, consider a model trained in 2022: if asked, “Where did the 2024 Olympic Games take place?”, the model would be unable to answer accurately, since that information did not exist in its training data.

To address such limitations, researchers developed the \textit{Retrieval-Augmented Generation} (RAG) framework, which enhances language models with access to external sources of information. RAG combines two key components: a retrieval step and a generation step.

    \subsection{Retrieving}
The retrieving step in a Retrieval-Augmented Generation (RAG) framework consists of fetching relevant information from a pre-built knowledge base, typically composed of a large collection of documents. But how is this database constructed?

The first step involves splitting the source documents into smaller segments, often called \textit{chunks}. This chunking process is essential, as working with smaller units improves both retrieval efficiency and semantic precision. For example, retrieving a short paragraph is usually more useful than retrieving an entire document, especially when answering a specific question.

Once the documents are chunked, the second step is to convert each chunk into a vector representation—a process known as embedding. These vectors encode the semantic meaning of each chunk and allow the retriever to compare a user's query with the stored knowledge in a mathematically meaningful way. Depending on the type of embeddings used, different retriever paradigms can be applied:
\begin{itemize}
    \item \textbf{Sparse retrievers}: These systems use \textit{sparse embeddings}, where most vector components are zero. One of the most well-known sparse retrieval methods is BM25, based on term frequency-inverse document frequency (TF-IDF) weighting. In this approach, words are represented by high-dimensional vectors, where each dimension corresponds to a term in the vocabulary. For instance, a word is encoded by placing a “1” at its corresponding position in the vocabulary and “0” elsewhere. A sentence or chunk is then represented by combining these one-hot vectors, typically by summing or averaging them. Sparse embeddings are interpretable and easy to implement, as they do not require training. However, they lack semantic understanding: two semantically similar words (e.g., “car” and “vehicle”) will be treated as unrelated if they do not share surface forms. As such, sparse retrievers do not capture relationships or meanings between words, limiting their expressiveness.
    \item \textbf{Dense retrievers}: In contrast, dense retrieval methods embed queries and chunks into a continuous, low-dimensional vector space using deep neural networks (often pre-trained language models). These embeddings capture richer semantic relationships between words, phrases, and concepts. For example, even if a query uses the word “automobile” and a chunk contains the word “car,” a dense retriever might still match them correctly due to their proximity in the embedding space. Dense retrievers are typically trained using contrastive learning techniques, where pairs of related (query, document) and unrelated examples are used to optimize the embedding space. Popular dense retrievers include DPR (Dense Passage Retrieval), which uses a dual-encoder architecture to independently embed queries and documents.
\end{itemize}
In summary, the effectiveness of the retrieval step largely depends on how the documents are chunked, embedded, and indexed. Sparse retrievers are fast, interpretable, and require no training, but they are limited in semantic capability. Dense retrievers, while more computationally intensive and requiring training data, provide better performance on tasks that demand deeper understanding of language.

Now that the documents have been embedded and stored in the vector database, the RAG (Retrieval-Augmented Generation) workflow proceeds to the retrieval phase. When a user issues a query, the system first embeds the query into the same vector space as the documents. This ensures that both queries and documents can be compared meaningfully based on their semantic content.

To determine which documents are most relevant to the query, a similarity metric is used to compare the query vector with the document vectors. One of the most commonly used metrics in this context is the \textbf{cosine similarity}.

Cosine similarity measures the cosine of the angle between two vectors in a high-dimensional space. It is defined as:

\[
\text{cosine\_similarity}(\vec{q}, \vec{d}) = \frac{\vec{q} \cdot \vec{d}}{\|\vec{q}\| \|\vec{d}\|}
\]

where:
\begin{itemize}
    \item $\vec{q}$ is the embedding of the query,
    \item $\vec{d}$ is the embedding of a document chunk,
    \item $\cdot$ denotes the dot product between vectors,
    \item and $\|\cdot\|$ represents the Euclidean norm (vector magnitude).
\end{itemize}

The cosine similarity score ranges from $-1$ to $1$, but in practice, since embeddings from modern language models are typically in the positive orthant, the values are between $0$ and $1$. A score close to $1$ indicates that the vectors are pointing in similar directions, meaning the query and the document are semantically similar. A score near $0$ suggests little to no similarity.

This metric is particularly effective because it focuses on the \textit{orientation} rather than the \textit{magnitude} of the vectors, making it well-suited for comparing text embeddings where the scale may vary but the direction encodes semantic meaning.

    \subsection{Generation}
The generation step refers to the phase in the RAG (Retrieval-Augmented Generation) pipeline where the retrieved documents are integrated into the prompt that is subsequently processed by the language model. This step is crucial, as the way information is presented to the model has a direct impact on the quality and relevance of the generated output.

When dealing with parameter-inaccessible generators — that is, black-box language models whose internal weights cannot be modified or fine-tuned — the most common and straightforward strategy for integration is known as \textbf{input-layer integration}. In this method, the retrieved documents are injected directly into the input of the model by concatenating them with the original prompt. This approach is simple to implement and compatible with any model that accepts plain text as input, making it widely adopted in practice.

A representative example of this strategy can be found in the work of \cite{inContextR}, where the retrieved documents are concatenated with the query or instruction that constitutes the base prompt. The resulting composite prompt is then fed into the language model in a single forward pass.

However, input-layer integration has a major drawback: it is subject to the maximum token limit of the model’s input space. Since most large language models have a fixed upper bound on the number of tokens they can process in one pass (for instance, 2048, 4096, or 8192 tokens depending on the model), this constraint significantly limits the number of documents that can be included in the prompt. As a result, potentially useful information may be excluded, which can degrade the performance on complex or information-rich tasks.

To address this limitation, one commonly used strategy is to summarize the retrieved documents before inserting them into the prompt. By applying either extractive or abstractive summarization techniques, it is possible to reduce each document to its most salient points, thereby saving tokens while preserving the essential information. This summarization step allows for a greater number of documents — or a broader scope of knowledge — to be incorporated into the prompt, all while staying within the token budget of the language model. Consequently, the model is exposed to a more comprehensive context, which can lead to more accurate and informative generation.

\section{Practical Implementation}
    \subsection{LangChain}
    \subsection{Ollama}
    \subsection{ChromaDB for Retrieving}
    \subsection{Generation}

\section{Evalutation}
    \subsection{Metrics}
    \subsection{Results}
    \subsection{Interpretation}

    
    
    
    


    